\item \textbf{Problema de repaso.} El techo perfectamente plano de una casa forma un ángulo $\theta$ con la horizontal. Cuando cada día su temperatura cambia, entre $T_c$ antes del amanecer y $T_h$ a media tarde, el techo se expande y contrae uniformemente con un coeficiente de expansión térmica $\alpha_1$. Sobre el techo hay una placa metálica rectangular plana con coeficiente de expansión $\alpha_2$ mayor que $\alpha_1$. La longitud de la placa es $L$, medida a lo largo de la pendiente del techo. La componente del peso de la placa perpendicular al techo se soporta mediante una fuerza normal distribuida uniformemente sobre el área de la placa. El coeficiente de fricción cinética entre la placa y el techo es $\mu_k$. La placa siempre está a la misma temperatura que el techo; así, suponga que su temperatura cambia de manera continua. Debido a la diferencia en coeficientes de expansión, cada trozo de la placa se mueve en relación con el techo bajo ella, excepto por puntos a lo largo de cierta línea horizontal que corre a través de la placa llamada línea estacionaria. Si la temperatura se eleva, partes de la placa bajo la línea fija se mueven abajo en relación con el techo y sienten una fuerza de fricción cinética que actúa hacia arriba en el techo. Elementos del área arriba de la línea fija se deslizan hacia arriba del techo y sobre ellos la fricción cinética actúa hacia abajo, paralela al techo. La línea fija no ocupa área, así, suponga que sobre la placa no actúa fuerza de fricción estática mientras la temperatura cambia. La placa como un todo casi está en equilibrio, así que la fuerza de fricción neta sobre ella debe ser igual a la componente de su peso que actúa hacia abajo del plano inclinado. 
(a) Demuestre que la línea fija está a una distancia de
$$\frac{L}{2} \left( 1 - \frac{\tan \theta}{\mu_s} \right)$$
abajo del borde superior de la placa. 
(b) Analice las fuerzas que actúan en la placa cuando la temperatura cae y demuestre que la línea estacionaria está a esa misma distancia sobre el borde inferior de la placa. 
(c) Demuestre que la placa baja del techo como un gusano moviéndose cada día una distancia
$$\frac{L}{\mu_k} (\alpha_2 - \alpha_1)(T_h - T_c) \tan \theta$$
(d) Evalúe la distancia que una placa de aluminio se mueve cada día si su longitud es de 1.20 m, la temperatura va en ciclos entre 4.00 y 36.0 °C y el techo tiene una pendiente de 18.5°, $1.50 \times 10^{-5} \text{ (}^\circ\text{C)}^{-1}$ de coeficiente de expansión lineal y 0.420 de coeficiente de fricción con la placa. 
(e) \textbf{¿Qué pasaría si?} ¿Y si el coeficiente de expansión de la placa es menor que el del techo? ¿La placa sube lentamente por el techo?
